%% ============================================================================
%%  front/abstract_fa.tex
%%  چکیده پایان‌نامه به زبان فارسی و کلیدواژه‌ها
%%  طراحی شده برای قالب رسمی دانشگاه خوارزمی
%% ============================================================================
%%
%%  توضیح:
%%  بر اساس شیوه‌نامه نگارش دانشگاه خوارزمی:
%%  - متن چکیده باید شامل بیان هدف، روش پژوهش، و یافته‌های کلیدی باشد.
%%  - تعداد کلمات چکیده معمولاً بین ۲۰۰ تا ۳۰۰ کلمه است.
%%  - در پایان چکیده، کمینه ۳ و بیشینه ۵ کلیدواژه قید می‌گردد.
%% ============================================================================

\abstractfa{%
امروزه تدوین و صفحه‌آرایی مستندات و رساله‌های علمی با رعایت دقیق استانداردهای نگارشی دانشگاهی نقشی اساسی در ارتقای کیفیت آثار پژوهشی ایفا می‌کند. در این پژوهش، یک سامانه جامع و استاندارد برای حروف‌چینی پایان‌نامه‌ها و رساله‌های دانشگاه خوارزمی بر پایه سامانه \lr{\XeLaTeX} و بسته زی‌پرشین طراحی و پیاده‌سازی شده است. در این ساختار، کلیه ضوابط مصوب دانشگاه شامل فواصل حاشیه‌ها (بالا ۳/۵، پایین ۲/۵، راست ۳/۰ و چپ ۳/۰ سانتی‌متر)، قلم‌های رسمی بی زر و بی تیتر، اندازه و فواصل دقیق سرعنوان‌ها، شماره‌گذاری فارسی تمامی عناصر (بخش‌ها، جدول‌ها، تصویرها، نمودارها و روابط ریاضی)، سرصفحه‌های متقارن، پانویس‌های استاندارد هر صفحه، و تفکیک منظم اجزای پیش از متن و پس از متن به‌طور دقیق رعایت شده است. علاوه‌بر این، سازوکار مدیریت مراجع دوزبانه (فارسی و انگلیسی) با استفاده از شیوه‌نامه‌های استاندارد بیب‌تک پیاده‌سازی شده تا ارجاع‌دهی به مقالات و کتاب‌ها بدون تداخل در جهت متن انجام پذیرد. نتایج ارزیابی‌های فنی نشان می‌دهد که قالب حاصل ضمن پایداری کامل در فرآیند کامپایل و فقدان هرگونه نویسه گم‌شده، کاربری آسانی برای دانشجویان تحصیلات تکمیلی به همراه دارد.
}

\keywords{قالب پایان‌نامه لاتک، زی‌پرشین، دانشگاه خوارزمی، مراجع دوزبانه، استانداردهای نگارشی.}

\makeAbstractFa
